\documentclass[11pt,a4paper]{article}

\usepackage[margin=2.5cm]{geometry}
\usepackage{amsmath,amssymb,booktabs}
\usepackage[T1]{fontenc}
\usepackage{lmodern}
\usepackage{hyperref}

\title{Pump-Current Conversions in the \texttt{twpa\_jax} Maps}
\author{Thesis working note}
\date{18 July 2026}

\begin{document}
\maketitle

\section{Purpose}

The gain-map programs convert an external pump-power coordinate into an
on-chip current before solving the nonlinear IPM circuit.  Two different
conventions have been used historically through the option
\texttt{--pump-current-jc-scale}.  They must not be confused with the ordinary
RMS-to-peak or Fourier-coefficient conversions.

\section{Power to physical peak current}

For an external pump power $P_{\mathrm{dBm}}$, after subtracting the line
attenuation $A_{\mathrm{dB}}$, the on-chip power is
\begin{equation}
  P_{\mathrm{W}}
  =10^{-3}10^{(P_{\mathrm{dBm}}-A_{\mathrm{dB}})/10}.
\end{equation}
The map code then computes
\begin{equation}
  I_{\mathrm{peak}}
  =\sqrt{\frac{2P_{\mathrm{W}}}{Z_0}}.
  \label{eq:peak-current}
\end{equation}
Thus the current stored in each map point is already a peak current.  In
particular, multiplying it by $\sqrt{2}$ again would be incorrect.

\section{Fourier and phasor bookkeeping}

For a real single-tone current, the conventions are
\begin{align}
  i(t) &= I_{\mathrm{peak}}\cos(\omega t+\phi),\\
       &= 2\operatorname{Re}\!\left[I_+e^{i\omega t}\right],\\
  I_+ &= \frac{I_{\mathrm{peak}}e^{i\phi}}{2}.
\end{align}
Equivalently, with an RMS RF phasor $I_{\mathrm{rms}}$,
\begin{equation}
  i(t)=\sqrt{2}\operatorname{Re}\!\left[I_{\mathrm{rms}}e^{i\omega t}\right],
  \qquad I_{\mathrm{rms}}=\frac{I_{\mathrm{peak}}}{\sqrt{2}},
  \qquad I_+=\frac{I_{\mathrm{rms}}}{\sqrt{2}}.
\end{equation}
The positive-phasor pump basis reconstructs the real waveform as
$2\operatorname{Re}\sum_k X_k e^{ik\omega_pt}$.  This factor of two belongs
to the reconstruction of a positive-frequency coefficient; it is not an
instruction to multiply the already-computed map current by two.

\section{The two map conversions}

Both paths use Eq.~\eqref{eq:peak-current}.  They differ only in the extra
source multiplier applied before the pump solve:
\begin{equation}
  I_{\mathrm{solve}}=s_{\mathrm{JC}}I_{\mathrm{peak}},
  \qquad s_{\mathrm{JC}}=\texttt{pump-current-jc-scale}.
\end{equation}

\begin{table}[h]
\centering
\begin{tabular}{@{}llll@{}}
\toprule
Setting & Source current & Meaning & Use \\
\midrule
$s_{\mathrm{JC}}=2.0$
  & $2I_{\mathrm{peak}}$
  & Historical JC-parity calibration
  & Legacy comparisons only \\
$s_{\mathrm{JC}}=1.0$
  & $I_{\mathrm{peak}}$
  & Current validated Python/IPM conversion
  & Current map regeneration \\
\bottomrule
\end{tabular}
\caption{The two source-current conventions.}
\end{table}

The factor-two path doubles the current amplitude.  If interpreted as a
matched-port power, that corresponds to a factor of four in power, or
$10\log_{10}(4)\simeq6.02$ dB.  It therefore moves the nonlinear operating
point substantially and can make a gain ridge appear at an apparently
unreachable external power.

The factor-one path does not add this historical JC multiplier.  As of
18~July~2026 it is the validated conversion for regenerating the Python/IPM
reference maps, and should be selected explicitly with
\begin{verbatim}
--pump-current-jc-scale 1.0
\end{verbatim}
The command-line default remains $2.0$ for backwards compatibility; relying
on the default is consequently unsafe for new map regeneration.

\section{Practical rule}

Use the following checklist:
\begin{enumerate}
  \item Convert dBm to $I_{\mathrm{peak}}$ using Eq.~\eqref{eq:peak-current}.
  \item Do not apply an additional RMS/peak factor.
  \item Use \texttt{--pump-current-jc-scale 1.0} for current Python/IPM maps.
  \item Use \texttt{2.0} only when deliberately reproducing a historical
        JC-parity run.
\end{enumerate}

\section*{Addendum (5 August 2026): Norton port termination}

Equation~\eqref{eq:peak-current} above is the \emph{travelling-wave} relation
between a wave-amplitude current and its power, $P=\tfrac12 I_{\mathrm{peak}}^2Z_0$.
It is not the relation that applies to this circuit. Every port of
\texttt{designs/ipm\_2c\_fixed} has a $G$-matrix diagonal entry of exactly
$0.02\,\mathrm S=1/Z_0$ and nothing else conductive, i.e.\ each drive is an
ideal current source in parallel with a $Z_0$ termination -- a
\textbf{Norton} source. The load only sees $I_{\mathrm{peak}}/2$, so the
available power is
\begin{equation}
  P_{\mathrm W}=\tfrac18 I_{\mathrm{peak}}^2 Z_0,
  \label{eq:peak-current-norton}
\end{equation}
a factor of 4 below the travelling-wave relation used to derive
Eq.~\eqref{eq:peak-current}, i.e.\ the inverse
conversion for a fixed on-chip power target needs
\emph{twice} the peak current this document's Eq.~\eqref{eq:peak-current}
computes:
\begin{equation}
  I_{\mathrm{peak}}=\sqrt{\frac{8P_{\mathrm W}}{Z_0}}
  \qquad\text{(Norton, not \eqref{eq:peak-current})}.
\end{equation}
Confirmed against the solver's own outgoing-power observable
(\texttt{pump\_outgoing\_power\_w}): at $I=7.2311\times10^{-6}\,\mathrm A$ it
reads $-64.857$~dBm, matching $I^2Z_0/8$ to the reported digits and
$6.0206$~dB ($=10\log_{10}4$) below $I^2Z_0/2$.

\textbf{This is orthogonal to \texttt{--pump-current-jc-scale}.} That
parameter is a separate multiplier on the source current chosen for
JosephsonCircuits.jl parity, layered on top of whichever power convention is
in effect -- it says nothing about whether the port itself is Norton or
travelling-wave, and this addendum does not change its recommended value.
The current code (\texttt{src/twpa\_solver/ports.py},
\texttt{--power-convention}, default \texttt{norton}) implements
Eq.~\eqref{eq:peak-current-norton}; \texttt{--power-convention
legacy\_traveling\_wave} reproduces every number this document's original
sections describe, bit for bit, for reproducing historical runs. See
\texttt{docs/development/psat\_comparison\_fix\_plan.md} Phase~1 and
\texttt{CLAUDE.md}'s "Port power convention: Norton, not travelling-wave" for
the full derivation and the depletion cross-check that confirmed it.

\end{document}
